

The background model in this analysis consists of nine components, grouped according to their spectra in the ROI or the uncertainty on their rate. Table~\ref{tab:bkgds} lists the expected number of events from each component.

\begin{table}[htbp]
    \caption{Number of events from various sources in the \SI{60}{d}$\times$\SI{5.5}{t} exposure~\cite{aalbers2022background}. The middle column shows the predicted number of events with uncertainties as described in the text. The uncertainties are used as constraint terms in a combined fit of the background model plus a \SI{30}{GeV/c\squared} WIMP signal to the selected data, the result of which is shown in the right column. $^{37}$Ar and detector neutrons have non-Gaussian prior constraints and are totaled separately. Values at zero have no lower uncertainty due to the physical boundary.}
    \label{tab:bkgds}
    \centering
    \begin{tabular}{>{\centering}p{3.3cm}>{\centering}p{2.38cm}>{\centering}p{2.3cm}}
    \tabularnewline
    \hline
    \hline
    Source & Expected Events & Fit Result\tabularnewline
    \hline
    $\beta$ decays + Det. ER &\centering 215 $\pm$ 36\phantom{0} &\centering 222 $\pm$ 16\phantom{0}  \tabularnewline
    $\nu$ ER& 27.1 $\pm$ 1.6\phantom{0} & 27.2 $\pm$ 1.6\phantom{0} \tabularnewline
    \XeOneTwoSeven & 9.2 $\pm$ 0.8 & 9.3 $\pm$ 0.8 \tabularnewline 
    \ce{^{124}Xe}& 5.0 $\pm$ 1.4 & 5.2 $\pm$ 1.4 \tabularnewline
    \ce{^{136}Xe} &	15.1 $\pm$ 2.4\phantom{0} & 15.2 $\pm$ 2.4\phantom{0} \tabularnewline
    $^8$B CE$\nu$NS & 0.14 $\pm$ 0.01 & 0.15 $\pm$ 0.01 \tabularnewline
    Accidentals	& 1.2 $\pm$ 0.3 & 1.2 $\pm$ 0.3 \tabularnewline
    \hline
    \rule{0pt}{1.0ex} Subtotal & 273 $\pm$ 36\phantom{0} 
    & 280 $\pm$ 16\phantom{0} \tabularnewline
    \hline \tabularnewline[-2.2ex]
    \ce{^{37}Ar} & [0, 288] & $52.5^{+9.6}_{-8.9}$\phantom{0} \tabularnewline[0.25ex]
    Detector neutrons &	$0.0^{+0.2 }$ & $0.0^{+0.2}$ \tabularnewline[0.25ex]
    \SI{30}{GeV/c\squared} WIMP &\centering -- &\centering $0.0^{+0.6}$ \tabularnewline[0.25ex]
    \hline
    Total & -- & 333 $\pm$ 17\phantom{0} \tabularnewline
    \hline
    \hline
    \end{tabular}\\
\end{table}



The dominant ER signal in the search comes from radioactive decay of impurities dispersed in the xenon. $^{214}$Pb from the $^{222}$Rn decay chain, $^{212}$Pb from $^{220}$Rn, and $^{85}$Kr have broad energy spectra that are nearly flat in energy across the ROI and are summed into an overall $\beta$ background. The concentrations of $^{214}$Pb (\SI{3.26}{\micro Bq/kg}) and $^{212}$Pb (\SI{0.14}{\micro Bq/kg}) are determined by fitting to energy peaks outside the ROI. 
The xenon was purified of krypton above ground using gas chromatography~\cite{akerib2018chromatography}, and an \textit{in situ} mass spectroscopy measurement of \SI{144(22)}{ppq} $^{\mathrm{nat}}$Kr~(g/g) informs the $^{85}$Kr rate estimate. 
The $\beta$ component is further combined with a small ($<$\SI{1}{\percent}) and similarly flat ER contribution from $\gamma$ rays originating in the detector components~\cite{akerib2020cleanliness} and cavern walls~\cite{AKERIB2020102391}. Solar neutrinos are also predicted to contribute a nearly flat ER spectrum in the ROI, with a rate calculated using Refs.~\cite{baxter2021recommended, Agostini2019Simultaneous, vinyoles2017new, aharmin2013combined}. As the prediction is very precise, neutrinos are kept separate from the detector $\beta$ background in this model. The naturally occurring isotopes of $^{124}$Xe (double electron capture) and $^{136}$Xe (double $\beta$ decay) contribute ER events, and the predictions are driven by the known isotopic abundances, lifetimes, and decay schemes~\cite{BerglundWieserIso,aprile2019observation,PhysRevC.89.015502}. 

Cosmogenic activation of the xenon prior to underground deployment produces short-lived isotopes that decayed during this first run, notably $^{127}$Xe (\SI{36.3}{d}) and $^{37}$Ar (\SI{35.0}{d})~\cite{TabRad_v7,TabRad_v8, PhysRevD.96.112011}. Atomic de-excitations following $^{127}$Xe L- or M-shell electron captures fall within the ROI if the ensuing $^{127}$I nuclear de-excitation $\gamma$ ray(s) escapes the TPC.  The rate of $^{127}$Xe electron captures is constrained by the rate of K-shell atomic de-excitations, which are outside the ROI.  The skin is effective at tagging the $^{127}$I nuclear de-excitation $\gamma$ ray(s), reducing this background by a factor of 5. 
The number of $^{37}$Ar events is estimated by calculating the exposure of the xenon to cosmic rays before it was brought underground, then correcting for the decay time before the search~\cite{aalbers2022cosmogenic}. A flat constraint of 0 to three times the estimate of 96 events is imposed because of large uncertainties on the prediction.

The NR background has contributions from radiogenic neutrons and coherent elastic neutrino-nucleus scattering (CE$\nu$NS) from $^8$B solar neutrinos. The prediction for the CE$\nu$NS rate, calculated as in Refs.~\cite{baxter2021recommended, Agostini2019Simultaneous, vinyoles2017new, aharmin2013combined}, is small due to the S2$>$600 phd requirement. The rate of radiogenic neutrons in the ROI is constrained using the distribution of single scatters in the FV tagged by the OD and then applying the measured neutron tagging efficiency from the AmLi calibration sources (\SI{89(3)}{\percent}). A likelihood fit of the NR component in the OD-tagged data is consistent with observing zero events, leading to a data-driven constraint of $0.0^{+0.2}$ applied to the search. This rate agrees with simulations based on detector material radioassay~\cite{akerib2020cleanliness}.


Finally, the expected distribution of accidentals is determined by generating composite single-scatter event waveforms from isolated S1 and S2 pulses and applying the WIMP analysis selections. The selection efficiency is then applied to unphysical drift time single-scatter-like events to constrain the accidentals rate.
